\chapter{Contexte général du projet}
\label{chap:contexte}

\begin{introchapitre}
Ce chapitre situe le projet. Il présente d'abord l'organisme d'accueil, puis le cadre et la
problématique du stage. Il propose ensuite une étude des solutions de supervision existantes,
dont les limites motivent la solution proposée, et se termine par la méthodologie de travail
et le planning suivis.
\end{introchapitre}

\section{Présentation de l'organisme d'accueil}

\subsection{\entreprise}

Ce stage s'est déroulé au sein de \textbf{\entreprise}, \afaire{présentation de l'entreprise :
date de création, siège, secteur d'activité, taille de l'équipe, principaux clients ou
produits (3 à 5 lignes)}.

\subsection{Activités et services}

\afaire{Décrire les principaux domaines d'activité de l'entreprise, par exemple :
développement de solutions logicielles, services cloud, intégration, conseil. Une liste à
puces de 3 à 6 éléments suffit.}

\subsection{Organisation}

\afaire{Présenter la structure de l'entreprise et situer l'équipe ou le département dans
lequel le stage a été effectué, sous la responsabilité de \encadrantPro.}

\capture[0.8\linewidth]{cap-01-organigramme.png}{Organigramme de \entreprise}{fig:organigramme}

\section{Cadre du projet}

Ce travail a été réalisé dans le cadre du stage du \diplome\ à \ecole, au sein de
\entreprise, \periodeStage. Il porte sur la conception et la réalisation complète d'un produit
logiciel, de l'étude des besoins jusqu'au déploiement, en suivant les pratiques DevOps et
DevSecOps : conteneurisation, intégration continue, tests automatisés et sécurité intégrée
dès la conception.

\section{Problématique}

Superviser une infrastructure consiste à répondre en permanence à trois questions :
\emph{quelque chose ne va-t-il pas ?}, \emph{pourquoi ?} et \emph{qui doit agir ?}
Les équipes réduites (startups, petites entreprises, projets étudiants) rencontrent trois
difficultés :

\begin{enumerate}
  \item \textbf{La complexité de mise en place.} La pile de référence associe Prometheus
        (collecte), Alertmanager (routage des alertes) et Grafana (visualisation). Chaque outil
        a sa configuration, son format et son cycle de mise à jour ; les relier demande du
        temps et des compétences que ces équipes n'ont pas toujours.
  \item \textbf{Le bruit des alertes.} Des alertes fondées uniquement sur des seuils fixes se
        déclenchent sur des pics passagers, se répètent pour un même problème et ne sont pas
        regroupées, ce qui provoque une \og fatigue d'alerte \fg{} : les vraies alertes finissent
        par être ignorées \cite{googlesre}.
  \item \textbf{Le manque d'explication.} Une alerte indique qu'un seuil est dépassé, mais
        pas quel processus en est la cause, ni depuis quand la situation dure.
\end{enumerate}

\begin{aretenir}[Problématique]
Comment offrir à une petite équipe une supervision \textbf{installable en une commande}, qui
produit \textbf{peu d'alertes mais des alertes fiables et expliquées}, tout en restant
\textbf{sécurisée} et \textbf{maintenable} comme un produit professionnel ?
\end{aretenir}

\section{Étude de l'existant}

\subsection{Notions de base}

\begin{itemize}
  \item \textbf{Métrique} : valeur numérique mesurée à intervalle régulier (utilisation CPU,
        mémoire, nombre de requêtes). C'est la matière première de la supervision.
  \item \textbf{Règle d'alerte} : condition sur une métrique, par exemple
        \texttt{cpu > 85~\%} pendant 60 secondes.
  \item \textbf{Incident} : ensemble d'alertes liées à une même cause.
  \item \textbf{Observabilité} : capacité à comprendre l'état interne d'un système à partir
        de ses métriques, de ses journaux et de ses traces. La supervision en est le premier
        niveau.
\end{itemize}

\subsection{Solutions existantes}

\paragraph{Prometheus, Alertmanager et Grafana.}
Prometheus collecte des séries temporelles en interrogeant des \emph{exportateurs} et évalue
des règles d'alerte écrites en PromQL, avec une clause \code{for} qui impose une durée avant
le déclenchement \cite{prometheus-rules}. Alertmanager déduplique, regroupe et route les
alertes vers les canaux de notification \cite{alertmanager}. Grafana assure la
visualisation \cite{grafana}. Cette pile est un standard de l'industrie, mais elle se compose
de trois produits à configurer séparément et ne propose pas de détection d'anomalies
statistique native.

\paragraph{Zabbix.}
Solution de supervision complète et ancienne, avec agent, serveur, base de données et
interface web \cite{zabbix}. Très riche fonctionnellement, elle demande cependant une base de
données externe et une configuration conséquente.

\paragraph{Netdata.}
Agent de supervision temps réel très léger, avec des tableaux de bord prêts à l'emploi et une
détection d'anomalies intégrée \cite{netdata}. Il est orienté vers la visualisation fine d'un
serveur plutôt que vers la gestion du cycle de vie des alertes et des incidents.

\paragraph{Datadog.}
Plateforme d'observabilité en mode SaaS, très complète (métriques, journaux, traces,
détection d'anomalies) \cite{datadog}. Les données quittent l'infrastructure de l'entreprise
et la facturation dépend du nombre d'hôtes, ce qui peut être coûteux pour une petite équipe.

\subsection{Synthèse comparative}

\begin{table}[htbp]
\centering
\small
\begin{tabularx}{\linewidth}{L{3.3cm}YYYYY}
\toprule
\textbf{Critère} & \textbf{Prometheus + Grafana} & \textbf{Zabbix} & \textbf{Netdata} & \textbf{Datadog} & \textbf{SentinelOps} \\
\midrule
Auto-hébergé & oui & oui & oui & non (SaaS) & oui \\
Installation & 3 outils à relier & serveur + base externe & simple & agent + compte & une commande \\
Règles avec durée & oui (\code{for}) & oui & oui & oui & oui \\
Détection d'anomalies & non native & fonctions de tendance & oui & oui & oui (score z robuste) \\
Alerte expliquée (preuves, processus) & non & partiel & partiel & oui & oui \\
Regroupement en incidents & via Alertmanager & oui & limité & oui & oui \\
Coût & gratuit & gratuit & gratuit / payant & par hôte & gratuit \\
\bottomrule
\end{tabularx}
\caption{Comparaison des solutions de supervision étudiées}
\label{tab:comparatif}
\end{table}

\subsection{Critique de l'existant}

Les solutions open source sont puissantes mais demandent un effort d'intégration important,
tandis que la solution SaaS la plus complète a un coût récurrent et externalise les données.
Aucune ne réunit à la fois une installation en une commande, une gestion complète du cycle de
vie des alertes, des alertes expliquées et une sécurité par défaut pour une petite équipe.

\section{Solution proposée}

SentinelOps est une plateforme auto-hébergée qui se lance avec la commande \code{sen run} et
qui :
\begin{itemize}
  \item \textbf{collecte} toutes les cinq secondes les métriques de l'hôte, des conteneurs
        Docker et de tout exportateur Prometheus ;
  \item \textbf{détecte} les incidents par des règles de seuil et par un détecteur
        statistique d'anomalies ;
  \item \textbf{décide} grâce à un automate d'états qui filtre les pics passagers, évite les
        doublons et regroupe les alertes en incidents ;
  \item \textbf{explique} chaque alerte avec les valeurs mesurées et les processus les plus
        consommateurs au moment du déclenchement ;
  \item \textbf{notifie} par Slack et e-mail avec garantie de livraison ;
  \item \textbf{affiche} l'ensemble en temps réel dans une interface web.
\end{itemize}

SentinelOps ne rejette pas l'écosystème existant : tout exportateur Prometheus peut être
interrogé, ce qui permet de superviser des applications tierces sans développement
supplémentaire.

\section{Méthodologie de travail}

Le projet a suivi une \textbf{approche itérative et incrémentale}, guidée par une feuille de
route versionnée. Chaque version livre un produit utilisable, puis est revue avant d'ajouter
de nouvelles fonctionnalités. Chaque itération suit le même cycle : analyse, conception,
réalisation, tests et revue.

\begin{table}[htbp]
\centering
\small
\begin{tabularx}{\linewidth}{L{3.4cm}Y}
\toprule
\textbf{Version} & \textbf{Contenu} \\
\midrule
v0.1 --- Core & collecte, règles, déduplication, corrélation, tableau de bord, Slack/e-mail, déploiement Docker Compose \\
v0.2 --- Intelligence & couche de détection extensible, détection d'anomalies, gestion d'incidents, durcissement \\
v0.3 --- Assistant IA & explication par un modèle de langage local (optionnel), suivi des changements \\
v1.0 --- Plateforme & prévisions, coûts cloud, modules Terraform, remédiation sous validation humaine \\
\bottomrule
\end{tabularx}
\caption{Feuille de route versionnée du produit}
\label{tab:roadmap}
\end{table}

Le stage couvre la version 0.1 complète et le début de la version 0.2. Une étape importante
a été une \textbf{revue de code} complète de la première version, qui a identifié seize
problèmes (fonctionnels, de sécurité et de qualité) ; leur correction a conduit à la refonte
décrite dans les chapitres 3 à 5.

\section{Planning}

La figure~\ref{fig:gantt} présente le déroulement prévisionnel du stage.

\begin{figure}[htbp]
\centering
\resizebox{!}{0.4\textheight}{% Planning prévisionnel (à ajuster aux dates réelles du stage)
\begin{ganttchart}[
    hgrid, vgrid={*1{draw=black!12}},
    x unit=0.82cm, y unit chart=0.62cm,
    title label font=\sffamily\scriptsize,
    bar label font=\sffamily\footnotesize,
    group label font=\sffamily\footnotesize\bfseries,
    milestone label font=\sffamily\footnotesize\itshape,
    bar/.append style={fill=bleuSentinel!35, draw=bleuSentinel},
    group/.append style={fill=bleuSentinel},
    milestone/.append style={fill=rougeAfaire, draw=rougeAfaire}
  ]{1}{16}
  \gantttitle{Semaines du stage}{16} \\
  \gantttitlelist{1,...,16}{1} \\
  \ganttgroup{Étude}{1}{3} \\
  \ganttbar{Étude de l'existant, besoins}{1}{2} \\
  \ganttbar{Conception initiale}{2}{3} \\
  \ganttgroup{Version 0.1 (Core)}{3}{8} \\
  \ganttbar{Agent : collecte, règles}{3}{5} \\
  \ganttbar{Dashboard : API, interface}{4}{7} \\
  \ganttbar{Notifications Slack / e-mail}{6}{8} \\
  \ganttmilestone{Prototype v0.1}{8} \\
  \ganttgroup{Version 0.2 (refonte)}{9}{13} \\
  \ganttbar{Revue de code, migrations, outbox}{9}{10} \\
  \ganttbar{Automate d'alertes, anomalies}{10}{11} \\
  \ganttbar{Sécurité, conteneurs durcis}{11}{12} \\
  \ganttbar{Refonte de l'interface}{12}{13} \\
  \ganttgroup{Industrialisation}{12}{16} \\
  \ganttbar{Tests, CI, Trivy, Dependabot}{12}{14} \\
  \ganttbar{Déploiement cloud / VPS}{14}{15} \\
  \ganttbar{Rédaction du rapport}{13}{16} \\
  \ganttmilestone{Soutenance}{16}
\end{ganttchart}
}
\caption{Diagramme de Gantt du stage \afaire{à ajuster aux dates réelles}}
\label{fig:gantt}
\end{figure}

\begin{conclusionchapitre}
Ce chapitre a présenté le contexte du stage et montré, à travers l'étude de l'existant, le
besoin d'une solution de supervision simple à déployer, aux alertes fiables et expliquées.
Le chapitre suivant formalise ces besoins.
\end{conclusionchapitre}
